\chapter{Validacion}

La validacion combina observacion visual, indicadores numericos y pruebas repetidas. No basta con que el robot se mueva: debe alcanzar el objetivo, evitar colisiones y mantener un comportamiento estable.

\begin{longtable}{P{0.28\textwidth}P{0.42\textwidth}P{0.20\textwidth}}
\toprule
\textbf{Prueba} & \textbf{Criterio} & \textbf{Estado esperado}\\
\midrule
\endhead
Atraccion a objetivo & El robot reduce distancia hasta entrar en tolerancia. & Cumplido\\
Parada cerca de meta & La velocidad se reduce sin atravesar el objetivo. & Cumplido\\
Obstaculo frontal & El robot gira y rodea el obstaculo sin contacto. & Cumplido\\
Obstaculos laterales & El robot mantiene distancia minima. & Cumplido\\
Integracion en celda & El robot no invade zonas de otros equipos. & Pendiente de captura final\\
\bottomrule
\end{longtable}

\section{Evidencias a incluir}

La version final debe incorporar capturas de CoppeliaSim para cada parte, nombres de scripts, parametros principales y una descripcion breve de los resultados observados. Si se exportan trayectorias, pueden añadirse graficas de distancia al objetivo y distancia minima a obstaculos.
